\section{Reading information from graphs}

A graph can contain a large amount of information, even when we do not know the formula of the function. Reading a graph means extracting that information carefully.

The aim is not to guess. The aim is to connect visible geometric features with statements about the function.

\subsection*{Reading values}

If a point \((a,b)\) lies on the graph of \(f\), then
\[
f(a)=b.
\]
The first coordinate is the input. The second coordinate is the output.

\begin{examplebox}
Suppose the graph of \(f\) contains the point \((3,-2)\). Then
\[
f(3)=-2.
\]
If the graph also contains \((-1,4)\), then
\[
f(-1)=4.
\]
\end{examplebox}

This is the most basic skill in reading graphs. It should become automatic: point on graph means input-output pair.

\subsection*{Zeros and signs}

A zero of a function is an input where the output is zero. On the graph, zeros are the points where the graph meets the horizontal axis.

If
\[
f(a)=0,
\]
then \((a,0)\) is on the graph.

The sign of \(f(x)\) can also be read from the graph. Where the graph lies above the horizontal axis, \(f(x)>0\). Where it lies below the horizontal axis, \(f(x)<0\).

\begin{examplebox}
Suppose a graph crosses the horizontal axis at \(x=-2\) and \(x=5\). Then
\[
f(-2)=0,
\qquad
f(5)=0.
\]
If the graph is above the axis between these points, then
\[
f(x)>0
\qquad\text{for } -2<x<5.
\]
If it is below the axis to the right of \(5\), then
\[
f(x)<0
\qquad\text{for } x>5
\]
on the visible part of the graph.
\end{examplebox}

A graph therefore answers not only numerical questions, but also qualitative questions about positivity and negativity.

\subsection*{Increasing and decreasing behavior}

A function is increasing on an interval if its values go upward as the input moves to the right. It is decreasing if its values go downward as the input moves to the right.

At this stage, we do not need a formal definition with quantifiers. The basic reading is visual: moving from left to right, does the graph rise or fall?

\begin{examplebox}
If a graph rises from left to right on the interval \([1,4]\), then the function is increasing on that interval. This means that larger inputs in that interval give larger function values.

If the graph falls from left to right on \([4,7]\), then the function is decreasing on that interval.
\end{examplebox}

Later, derivatives will give a powerful method for studying increasing and decreasing behavior. But the geometric idea should be understood first.

\subsection*{Maximum and minimum values}

A graph can also show where a function has high or low values. A local maximum is a point where the graph is higher than nearby points. A local minimum is a point where the graph is lower than nearby points.

\begin{examplebox}
Suppose a graph reaches a visible peak at the point \((2,5)\). Then near \(x=2\), the function has a local maximum value \(5\).

If the graph reaches a visible valley at \((-1,-3)\), then near \(x=-1\), the function has a local minimum value \(-3\).
\end{examplebox}

The input tells us where the maximum or minimum occurs. The output tells us what the maximum or minimum value is.

\subsection*{Graph reading is not formula guessing}

Sometimes a graph is given without a formula. This is not a weakness. A graph can still answer many questions.

From a graph we may read:
\begin{itemize}
\item values such as \(f(2)\),
\item zeros of the function,
\item intervals where the function is positive or negative,
\item intervals where the function is increasing or decreasing,
\item approximate maximum and minimum values,
\item the domain and range visible in the picture.
\end{itemize}

But we should be honest about precision. A graph may show approximate values, while a formula may give exact values. Both forms of information are useful, but they are not the same.

\begin{remarkbox}
Reading a graph means translating geometric features into mathematical statements. It does not always mean finding a formula.
\end{remarkbox}

\begin{summarybox}
When reading a graph, ask:
\begin{itemize}
\item What is the input coordinate?
\item What is the output coordinate?
\item Where is the graph equal to zero?
\item Where is the graph above or below the horizontal axis?
\item Where does the graph rise or fall as we move from left to right?
\item Where are visible peaks or valleys?
\item Am I reading exact information or approximate information?
\end{itemize}
\end{summarybox}

A graph is a language. Learning to read it is part of learning functions.